\chapter{Introduction}
\label{ch:intro}

\section{Purpose}

This document records the tools considered for the GENESYS Translation Tool,
which were adopted, which were not, and the reasoning in each case. It is written
so that a reader who was not present for the evaluation can understand why the
toolchain looks the way it does, and so that a later maintainer can tell which
constraints are fixed and which were judgement calls.

Each tool is described under three headings: what it does, why it is useful, and
whether it works for this project. Rejected tools are documented at the same
length as adopted ones, because the reason a reasonable option was turned down is
more informative than the reason an obvious one was taken.

\section{Evaluation Criteria}

Tools were judged against five criteria, in this order of priority.

\begin{enumerate}[nosep]
	\item \textbf{Access.} Whether the team can obtain and use it, given a
	      University Edition licence and no access to the vendor's source.
	\item \textbf{Fidelity.} Whether it exposes the structure the translation
	      needs, particularly relationships between entities, rather than flattened
	      values.
	\item \textbf{Automatability.} Whether it can be driven programmatically
	      rather than through a user interface.
	\item \textbf{Stability.} Whether it is a supported interface that will
	      survive a version change, rather than an implementation detail.
	\item \textbf{Cost to the schedule.} Whether the effort to adopt it is
	      proportionate to a two-semester project with three people.
\end{enumerate}

Criterion 1 eliminated more candidates than the other four combined, and it is
why several otherwise reasonable options appear in Chapter~\ref{ch:rejected}.

\section{Scope}

This document covers the tools that affect the design of the deliverable: the
interfaces the tool talks to, the libraries it is built from, and the runtimes it
targets. Supporting tools the team uses to produce the work are covered in less
detail, since those choices are reversible. It does not cover GENESYS itself,
which is a given rather than a choice, nor the client's development environment.
